\documentclass[11pt]{article}
\usepackage[margin=1in]{geometry}
\usepackage{amsmath,amssymb}
\usepackage{booktabs}
\usepackage{hyperref}

\title{Mapping responses to focal injections of bicuculline in the lateral parafacial region identifies core regions for maximal generation of active expiration}
\author{}
\date{}

\begin{document}
\maketitle

\begin{abstract}
The lateral parafacial area (pFL) is a crucial region involved in respiratory control, particularly in generating active expiration through an expiratory oscillatory network. Active expiration involves rhythmic abdominal (ABD) muscle contractions during late expiration and increases ventilation during elevated respiratory demands. The precise anatomical location of the expiratory oscillator along the ventral medulla's rostro-caudal axis is debated. We employed bicuculline injections at various pFL sites ($-0.2$ to $+0.8$\,mm from VIIc) to investigate the impact of GABAergic disinhibition on respiration. These injections consistently elicited ABD recruitment, but response strength varied along the rostro-caudal zone. The most robust and enduring changes in tidal volume, minute ventilation, and combined respiratory responses occurred at more rostral pFL locations ($+0.6$/$+0.8$\,mm from VIIc). Multivariate analysis of the respiratory cycle further differentiated between locations, revealing the core site for active expiration generation.
\end{abstract}

\section{Introduction}
Respiratory control requires precise coordination among brainstem nuclei. The preB\"otzinger Complex (preB\"otC) generates the inspiratory rhythm and pattern \citep{gray2001,smith1991,tan2008,wang2014}. The lateral parafacial region (pFL, also called pFRG) contains neurons responsible for generating active expiration via recruitment of expiratory muscles such as the abdominal (ABD) muscles \citep{debritto2017,huckstepp2015,pagliardini2011,pisanski2019}. The expiratory oscillator's location remains debated: some studies implicate the caudal tip of the facial nucleus (VIIc, $-0.2$ to $+0.5$\,mm) \citep{huckstepp2015,pisanski2019}, others more rostral coordinates ($+0.3$ to $+1.0$\,mm) \citep{silva2019}.

We used focal injections of bicuculline, a competitive antagonist of GABA-A receptors, along the rostrocaudal axis of the ventral brainstem to enhance excitability at target coordinates and reveal the most efficient site for generating active expiration. We combined traditional respiratory measures (tidal volume, minute ventilation, respiratory rate, oxygen metabolism, ABD amplitude/area) with a novel multidimensional cycle-by-cycle analysis that characterises the effect of bicuculline on the late expiratory, inspiratory, and post-inspiratory phases.

\section{Methods}
\subsection{Animals and surgery}
Thirty-five adult male Sprague Dawley rats ($340.4 \pm 13.2$\,g) were used. Experiments were approved by the Animal Care and Use Committee of the University of Alberta (AUP\#461). Rats were anaesthetised with isoflurane (5\% induction, 1--3\% maintenance) and converted to urethane (1.5--1.7\,g/kg IV). A tracheal cannula was connected to a flow head, and oxygen consumption was computed from the fractional O$_2$ concentration and tracheal flow rate per Depocas and Hart (1957). Paired EMG electrodes were inserted into the oblique ABD and diaphragm (DIA) muscles; bilateral vagotomy was performed; body temperature was maintained at $37\pm1^\circ$C.

\subsection{Bicuculline injections}
Rats received bilateral pressure injections of 200 nL of 200\,$\mu$M bicuculline methochloride + fluorobeads in HEPES at 100\,nL/min (30\,$\mu$m tip). Injections were placed at $-0.2$ ($n=5$), $+0.1$ ($n=7$), $+0.4$ ($n=5$), $+0.6$ ($n=6$), or $+0.8$ ($n=5$) mm from VIIc, at 2.8\,mm lateral. HEPES + fluorobead controls: $n=7$. Physiological responses were recorded for 20--25\,min.

\subsection{Histology and anatomy}
Brains were transcardially perfused, cryoprotected, and sectioned at 30\,$\mu$m. Immunostaining used anti-ChAT (goat, 1:800), anti-PHOX2B (mouse, 1:100), anti-cFos (rabbit, 1:500), anti-TH (chicken, 1:800). Sections were examined at 120\,$\mu$m spacing, spanning 400\,$\mu$m caudal to 1000\,$\mu$m rostral to VIIc.

\subsection{Signal processing and analyses}
EMG was amplified $\times 10{,}000$ and filtered 300\,Hz--1\,kHz; signals were sampled at 1\,kHz and rectified with a 0.08\,s time constant. Airflow was integrated to obtain tidal volume $V_T$ using a five-point calibration curve (0.5--5\,mL range). Data were analysed in 2-minute bins over 20\,min post-injection; mean-cycle measurements used $\sim$300 cycles per baseline bin and $\sim$100 cycles per response bin. Each respiratory cycle was split into late-E, inspiratory, and post-I phases using zero-crossings of mean airflow. The area under the curve per phase was computed, with baseline AUC subtracted.

For the multidimensional cycle-by-cycle analysis, each timepoint within a respiratory cycle was placed as a point in a 3D space defined by airflow, $\int$DIA EMG, and $\int$ABD EMG. The straight-line (Euclidean) distance between response and baseline trajectories, and the Mahalanobis distance (scaled by baseline cycle-to-cycle covariance), were computed per phase. Statistical comparisons across groups used repeated-measures ANOVA with post-hoc comparisons.

\section{Results}
\subsection{Bicuculline elicits ABD recruitment across the rostrocaudal axis}
Bicuculline injections at every tested coordinate ($-0.2$ to $+0.8$\,mm from VIIc) recruited ABD activity and generated active expiration; control (HEPES + fluorobead) injections did not. Response magnitude varied along the rostrocaudal axis.

\subsection{Rostral sites yield the strongest tidal volume and minute ventilation changes}
The strongest and most enduring changes in tidal volume ($V_T$), minute ventilation ($V_E$), and combined respiratory responses occurred at the two most rostral locations ($+0.6$ and $+0.8$\,mm from VIIc). ABD activation lasted longest at these sites, and the swiftest onset of the ABD response (shortest ABD delay) was observed at $+0.6$\,mm.

\subsection{Multivariate cycle analysis differentiates rostral sites}
Multivariate phase-by-phase analysis distinguished the responses at $+0.6$ vs.\ $+0.8$\,mm. Injections at $+0.8$\,mm produced the strongest deformations of the respiratory trajectory during late expiration (late-E) and post-inspiration (post-I), whereas injections at $+0.6$\,mm produced stronger deformations during the inspiratory phase. The Euclidean and Mahalanobis distances, computed per respiratory phase, were consistent across the two measures.

\begin{table}[h]
\centering
\caption{Summary of response strength by injection location.}
\begin{tabular}{lc}
\toprule
Injection location (mm from VIIc) & $n$ \\
\midrule
$-0.2$ & 5 \\
$+0.1$ & 7 \\
$+0.4$ & 5 \\
$+0.6$ & 6 \\
$+0.8$ & 5 \\
CTRL (HEPES + fluorobead) & 7 \\
\bottomrule
\end{tabular}
\end{table}

\section{Discussion}
Bicuculline injections along the rostro-caudal axis of pFL produce active expiration at every tested coordinate from $-0.2$ to $+0.8$\,mm, but the core site for maximal respiratory activation is rostral ($+0.6$ to $+0.8$\,mm) rather than at VIIc. This aligns with Silva et al.\ \citep{silva2019} and extends prior work that had focused on more caudal coordinates \citep{huckstepp2015,pisanski2019}. The multivariate trajectory analysis (Euclidean + Mahalanobis distances across the inspiratory, late-E, and post-I phases) discriminated the two rostral sites, with the $+0.8$\,mm site emphasising expiratory deformations and $+0.6$\,mm emphasising inspiratory deformations. The trajectory framework and AUC decomposition enable tracking of all three respiratory signals simultaneously and should be broadly useful for studies of expiratory oscillator modulation, plasticity, and therapeutic targeting in respiratory disorders.

\bibliographystyle{plain}
\bibliography{references}

\end{document}
